\documentclass{article}
\usepackage{amsmath, amssymb}
\begin{document}
\title{Digest: Trivial Zeros of $p$-adic $L$-functions at Near Central Points}
\author{Denis Benois (Digest)}
\maketitle

\section{Introduction}
This document digests Denis Benois's paper on the trivial zeros of $p$-adic $L$-functions of modular forms. The paper proves a Mazur-Tate-Teitelbaum style formula for the derivative of $p$-adic $L$-functions at near central points.

\section{The Main Theorem}
Let $f$ be a newform of weight $k$ and character $\varepsilon$. For a $p$-adic $L$-function $L_{p, \alpha}(f, \omega^{k_0}, s)$ having a trivial zero at $s=k_0 \ge k/2$, the derivative at the near central point is given by:
\begin{equation}
L_{p, \alpha}'(f, \omega^{k_0}, k_0) = \ell_{\alpha}(f) \left(1 - \frac{\varepsilon(p)}{p}\right) \tilde{L}(f, k_0)
\end{equation}
where $\ell_{\alpha}(f)$ is the generalized $\ell$-invariant.

\section{$(\varphi, \Gamma)$-modules and the $\ell$-invariant}
The construction relies heavily on the theory of $(\varphi, \Gamma)$-modules over the Robba ring $\mathcal{R}_L$. The $\ell$-invariant $\ell(V, D)$ is defined via the Galois cohomology of the potentially semistable representation $V$ and its crystalline submodule $D$. 

\section{Conclusion}
This framework provides a unified proof of the trivial zero conjecture for elliptic modular forms, generalizing Greenberg's $\ell$-invariant to non-ordinary cases.
\end{document}
